\section{Partition and Search Details}
\label{sec:structural-details}
\subsection{Count Monotonicity and Membership}
\label{sec:partition-properties}
Treat each point as a distinct atom. For a trie subtree rooted at $v$, let
$b_v(T)$ be its emitted block count and $U_v(T)$ the residual point set it
returns to its parent. The synthetic root never emits. For $T_1<T_2$, we
prove the stronger induction invariant
\begin{equation}
 b_v(T_1)\ge b_v(T_2),\qquad
 b_v(T_1)=b_v(T_2)\Longrightarrow U_v(T_1)\supseteq U_v(T_2).
\end{equation}
Let $d=\sum_c(b_c(T_1)-b_c(T_2))$ over children, and let $e_i\in\{0,1\}$
indicate whether $v$ emits at threshold $T_i$. The block-count difference
is $d+e_1-e_2$. If $d\ge1$, it is nonnegative. If $d=0$, all child counts
are equal; their residual sets at $T_1$ contain those at $T_2$. After adding
the same terminal points at $v$, the pending set at $T_1$ therefore contains
the pending set at $T_2$. A high-threshold emission without a low-threshold
emission is impossible, proving count monotonicity.

If the final counts are equal, either $d=0$ and $e_1=e_2$, preserving residual
containment or emptying both sets, or $d=1,e_1=0,e_2=1$, making the
high-threshold residual empty. At the synthetic root both emission indicators
are zero, and the childwise invariant passes to their union. This proves
both assertions by postorder induction. Taking complements at the root shows
that equal global block counts imply containment of the represented point
sets in the opposite direction. Equation~\ref{eq:block-count-bound} follows
because each disjoint emitted set contains at least $T+1$ points for integer $T\ge0$.

\paragraph{Equal counts do not imply nested blocks.}
Let terminal weights be $a:2,ab:1,ac:1,abc:2$, with $a<b<c$ and $N=6$.
At $T=1$, the blocks are $abc:\{abc\}$ (mass 2) and
$a:\{a,ab,ac\}$ (mass 4). At $T=2$, they are
$ab:\{ab,abc\}$ (mass 3) and $a:\{a,ac\}$ (mass 3).
Both layers have two blocks and cover all six points, but the original
$a$ block is split between the new blocks.

\paragraph{Represented and authorized mass behave differently.}
For weights $a:1,ab:2,abc:3,ac:2$, thresholds 2 and 4 represent eight and
five points, respectively; block count falls from two to one. Conversely,
weights $a:1,ab:3,ac:3$ represent six points at threshold 2 and seven at
threshold 6, again with a two-to-one count change. Query-authorized mass can
fall even on a count plateau: with $a:1,ab:3$ and $Q=\{b\}$, threshold 2
emits authorized block $ab$ of mass 3, whereas threshold 3 emits only the
unauthorized block $a$ of mass 4. Thus count monotonicity enables structural
threshold selection but does not predict query authorization.

\subsection{A Finite-Budget Two-Overlay Example}
\label{sec:finite-budget-example}
Consider singleton groups with vectors $y_0,\ldots,y_8$ and label sets,
in order, $a,ab,abc,abd,ac$, followed by $acd,ace,acf,ad$.
Concatenation denotes a set under $a<b<c<d<e<f$. For thresholds $T_1=1$ and $T_2=4$, the independent
partitions produce the direct memberships in Table~\ref{tab:seed-trace}.
At level 1, the $a$ block excludes the emitted $ab$ and $ac$ descendants;
at level 2, $a$ owns all nine vectors. Query $Q=\{a\}$ certifies all four
blocks. Its prefix-frontier group $a$ contributes $y_0$, retagged to level 1,
which duplicates that block's seed and is removed during seed deduplication.

\begin{table}[H]
\centering\small
\caption{Illustrative ownership and seed selection, with $K=1$ and $L_s=2$.
Stored block entries and query distances are chosen for this example.}
\label{tab:seed-trace}
\begin{tabular}{clll}
\toprule
Level & Root & Direct members & Entry (distance)\\
\midrule
1 & $a$ & $y_0,y_8$ & $y_0\ (3)$\\
1 & $ab$ & $y_1,y_2,y_3$ & $y_1\ (1)$\\
1 & $ac$ & $y_4,y_5,y_6,y_7$ & $y_4\ (2)$\\
2 & $a$ & $y_0,\ldots,y_8$ & $y_1\ (1)$\\
\bottomrule
\end{tabular}
\end{table}

The two retained seeds are $(y_1,1)$ and $(y_1,2)$, distinct level states of
the same vector. The first expansion uses only level-1 edges. On offering
$y_2$ at distance $0.5$, before any other improving offer, the queue retains
$(y_2,1)$ and $(y_1,1)$ and evicts
$(y_1,2)$ before that coarse state expands. This trace shows how extra
representations can compete for a fixed budget even with isolated adjacency.
It also shows why coverage of the full frontier does not imply coverage
after initialization and queue eviction.
